\documentclass[11pt]{article}
\usepackage[margin=0.75in]{geometry}
\usepackage{amsmath,amssymb,siunitx}
\usepackage{parskip}
\usepackage{enumitem}

\title{PS160 --- Midterm 3 (Final) Equation Sheet}
\author{Optics, EM Waves, Interference, Diffraction (Modules 33--36) \\
\textit{Plus reminders of Midterm 1\,\&\,2 material that must be combined with optics}}
\date{}

\begin{document}
\maketitle

\noindent\textbf{Important:} The official equation sheet provided on exam day
covers \emph{only} Midterm 1 and Midterm 2 material (Modules 12, 14--20).
\textbf{Nothing} from Modules 33--36 (this sheet's main topic) appears on the
official sheet. Every formula on this page must therefore be memorized or
quickly derivable. See \texttt{equation\_gaps.md} for the full audit.

\section*{Module 33 --- Nature and Propagation of Light, EM Waves}

\paragraph{Speed of light, index, wavelength in medium.}
\[
c = \SI{3.00e8}{m/s} = \frac{1}{\sqrt{\mu_0\varepsilon_0}}
\qquad
n = \frac{c}{v}
\qquad
\lambda_{\text{med}} = \frac{\lambda_0}{n}
\]
Frequency does not change at an interface; speed and wavelength scale by $1/n$.

\paragraph{Plane wave relations.} $E/B = c$, $\vec E\perp\vec B\perp\vec k$.

\paragraph{Energy density.} $u = \varepsilon_0 E^2 = B^2/\mu_0$.

\paragraph{Poynting vector and intensity.}
\[
\vec S = \frac{1}{\mu_0}\vec E\times\vec B,
\qquad
I = \tfrac{1}{2}c\varepsilon_0 E_{\max}^2 = \frac{E_{\max}^2}{2\mu_0 c}
\]

\paragraph{Radiation pressure.} $p_{\text{abs}} = I/c$, $p_{\text{refl}} = 2I/c$.

\paragraph{Law of reflection.} $\theta_r = \theta_a$ (from the normal).

\paragraph{Law of refraction (Snell).} $n_a\sin\theta_a = n_b\sin\theta_b$.

\paragraph{Total internal reflection (only when $n_a > n_b$).}
\[
\sin\theta_c = \frac{n_b}{n_a}
\]

\paragraph{Polarization (Malus's law).}
\[
I = I_{\max}\cos^2\phi
\qquad\text{(unpolarized through ideal polarizer: }I = I_0/2)
\]

\paragraph{Brewster's angle (reflected light fully polarized).}
\[
\tan\theta_B = \frac{n_2}{n_1}
\]

\section*{Module 34 --- Geometric Optics: Mirrors and Lenses}

\paragraph{Mirror / thin-lens equation.}
\[
\frac{1}{s} + \frac{1}{s'} = \frac{1}{f}
\qquad
f_{\text{mirror}} = \frac{R}{2}
\]

\paragraph{Lateral magnification.}
\[
m = \frac{y'}{y} = -\frac{s'}{s}
\]

\paragraph{Lensmaker's equation (thin lens in air).}
\[
\frac{1}{f} = (n - 1)\!\left(\frac{1}{R_1} - \frac{1}{R_2}\right)
\]

\paragraph{Refraction at a single spherical surface.}
\[
\frac{n_1}{s} + \frac{n_2}{s'} = \frac{n_2 - n_1}{R}
\]

\paragraph{Sign conventions (incoming light from left).}
\begin{itemize}\itemsep0pt
  \item $s > 0$ if object is on the incoming side (left).
  \item $s' > 0$ if image is on the outgoing side (transmitted side for a lens; same side as object for a mirror).
  \item $R > 0$ if the center of curvature is on the outgoing side.
  \item $f > 0$ converging (convex lens, concave mirror); $f < 0$ diverging (concave lens, convex mirror).
\end{itemize}

\paragraph{Two-lens systems.}
\[
s_2 = d - s_1', \qquad m_{\text{tot}} = m_1 m_2
\]

\paragraph{Optical instruments (angular magnification).}
\[
M_{\text{magnifier}} = \frac{25\;\text{cm}}{f}
\qquad
M_{\text{telescope}} = -\frac{f_{\text{obj}}}{f_{\text{eye}}}
\]

\paragraph{Camera $f$-number, lens power, vision.}
\[
f\text{-number} = \frac{f}{D},
\qquad
P\;[\text{diopter}] = \frac{1}{f\;[\text{m}]}
\]
Myopia (near-sighted) corrected by diverging lens ($P<0$);
hyperopia (far-sighted) corrected by converging lens ($P>0$).

\paragraph{Apparent depth.}
\[
d_{\text{app}} = d_{\text{real}}\cdot\frac{n_{\text{obs}}}{n_{\text{med}}}
\]

\section*{Module 35 --- Interference}

\paragraph{Phase from path difference.}
\[
\Delta\phi = \frac{2\pi}{\lambda}\Delta r
\]

\paragraph{Two-source conditions (in phase, coherent).}
\[
\Delta r = m\lambda \;\text{(bright)},
\qquad
\Delta r = (m+\tfrac{1}{2})\lambda \;\text{(dark)}
\]

\paragraph{Young's double-slit.}
\[
d\sin\theta = m\lambda \;\text{(bright)},
\qquad
y_m \approx \frac{m\lambda R}{d},
\qquad
\text{fringe spacing} = \frac{\lambda R}{d}
\]
Width of central maximum (between $m=\pm 1$ dark fringes): $w = 2\lambda R/d$.

\paragraph{Two-source intensity.}
\[
I(\theta) = I_0\cos^2\!\left(\frac{\pi d\sin\theta}{\lambda}\right)
\]

\paragraph{Thin-film interference (normal incidence).}
\[
2 n t = m\lambda_0 \quad\text{or}\quad (m+\tfrac{1}{2})\lambda_0
\]
Choice flips with the count of \emph{hard} reflections (low-$n$ to high-$n$): zero or two hard reflections give the integer rule for constructive; one hard reflection swaps the rule.

\paragraph{Michelson interferometer.}
\[
2\,\Delta d = \Delta m\,\lambda
\]

\section*{Module 36 --- Diffraction}

\paragraph{Single-slit (Fraunhofer) minima.}
\[
a\sin\theta = m\lambda,\qquad m = \pm 1,\pm 2,\dots
\]
Width of central max: $w = 2\lambda R/a$.

\paragraph{Single-slit intensity.}
\[
I = I_0\!\left(\frac{\sin(\beta/2)}{\beta/2}\right)^{\!2},
\qquad
\beta = \frac{2\pi a\sin\theta}{\lambda}
\]

\paragraph{Two-slit pattern (with finite slit width).}
\[
I = I_0\cos^2\!\left(\frac{\pi d\sin\theta}{\lambda}\right)
\!\left(\frac{\sin(\beta/2)}{\beta/2}\right)^{\!2}
\]

\paragraph{Diffraction grating.}
\[
d\sin\theta = m\lambda,\qquad
R = \frac{\lambda}{\Delta\lambda} = N m
\]
Visible orders: $m_{\max} = \lfloor d/\lambda\rfloor$.

\paragraph{Circular aperture (Rayleigh criterion).}
\[
\sin\theta_R = 1.22\,\frac{\lambda}{D}
\]

\paragraph{X-ray diffraction (Bragg's law).}
\[
2d\sin\theta = m\lambda \qquad(\theta\ \text{measured from the planes})
\]

\section*{Reminders from Midterm 1\,\&\,2 (these \emph{are} on the official sheet)}

\paragraph{Pressure, Bernoulli, continuity (M12).}
$\rho = m/V$,
$p = F_\perp/A$,
$p(y) = p_0 - \rho g y$,
$\rho A v = \mathcal{C}$,
$p + \rho g y + \tfrac{1}{2}\rho v^2 = \mathcal{C}'$.

\paragraph{SHM \& waves (M14--15).}
$\ddot x + \omega^2 x = 0$,
$x(t) = A\cos(\omega t + \phi)$,
$\omega = \sqrt{k/m}$, $\sqrt{g/\ell}$, $\sqrt{mgd/I}$;
$y(x,t) = A\cos(kx - \omega t)$, $v = \lambda f = \omega/k$;
$v_{\text{string}}=\sqrt{F/\mu}$, $v_{\text{solid}}=\sqrt{Y/\rho}$, $v_{\text{fluid}}=\sqrt{B/\rho}$;
$f_n = nv/(2L)$ or $nv/(4L)$;
$\lambda_n = 2L/n$ or $4L/n$;
$P_{av} = \tfrac{1}{2}\sqrt{\mu F}\,\omega^2 A^2$.

\paragraph{Sound (M16).}
$I = \tfrac{1}{2}\sqrt{\rho B}\,\omega^2 A^2 = p_{\max}^2/(2\sqrt{\rho B})$,
$p_{\max} = BkA$,
$I = P/(4\pi r^2)$,
$\beta = (10\,\text{dB})\log_{10}(I/I_0)$ with $I_0 = \SI{e-12}{W/m^2}$,
$f_L = f_S(v+v_L)/(v+v_S)$,
$f_{\text{beat}} = f_a - f_b$.

\paragraph{Heat \& kinetic theory (M17--18).}
$\Delta L = \alpha L_0\Delta T$, $\Delta V = \beta V_0\Delta T$, $\beta = 3\alpha$;
$Q = mc\Delta T = nC\Delta T$, $Q = \pm mL$;
$H = kA\,\Delta T/L$, $H = A\sigma T^4$ (sheet has no $e$);
$pV = NkT = nRT$, $kN_A = R$;
$\tfrac{1}{2}m\langle v^2\rangle = \tfrac{3}{2}kT$, $v_{rms} = \sqrt{\langle v^2\rangle}$;
$C_V = (f/2)R$, $C_P = C_V + R$, $\gamma = C_P/C_V$, $C_V = 3R$ (solid).

\paragraph{Thermo cycles (M19--20).}
$dU = dQ - dW$,
$W = \int p\,dV$,
$W = Q = Q_H + Q_C$ for a cycle,
$\Delta S \ge 0$,
$\Delta S = \int dQ/T$,
$S = k\ln w$;
$e = W/Q_H$,
$K = Q_C/(-W)$,
$K = -Q_H/(-W)$ (heat pump heating mode);
Carnot: $\Delta S = Q_H/T_H + Q_C/T_C = 0$.

\section*{Things NOT on the sheet that you must memorize}

A non-exhaustive list (see \texttt{equation\_gaps.md} for the full inventory):
\begin{itemize}\itemsep0pt
  \item Temperature scale conversions $T_K = T_C + 273.15$, $T_F = \tfrac{9}{5}T_C + 32$.
  \item Carnot efficiency $e = 1 - T_C/T_H$ (derive from sheet's $\Delta S = 0$ and $e = W/Q_H$).
  \item Adiabatic process: $pV^\gamma = \text{const}$, $TV^{\gamma-1} = \text{const}$, $W = (p_1V_1 - p_2V_2)/(\gamma - 1)$.
  \item Isothermal work $W = nRT\ln(V_2/V_1)$.
  \item Useful entropy: $\Delta S_{\text{isoT}} = nR\ln(V_2/V_1)$, $\Delta S_{\text{heat}} = mc\ln(T_2/T_1)$, $\Delta S_{\text{phase}} = \pm mL/T$.
  \item Temperature-dependent sound speed $v(T) = (331\,\text{m/s})\sqrt{T/(273\,\text{K})}$.
  \item Stefan--Boltzmann with emissivity: $H = e A\sigma T^4$.
  \item All of M33--M36 (above): Snell, Brewster, Malus, Maxwell relations, thin lens, lensmaker, $d\sin\theta=m\lambda$, $a\sin\theta=m\lambda$, Bragg, Rayleigh, thin-film $2nt$, etc.
\end{itemize}

\section*{Constants}
\begin{tabular}{ll}
Speed of light & $c = \SI{2.998e8}{m/s}$ \\
Permittivity of free space & $\varepsilon_0 = \SI{8.854e-12}{F/m}$ \\
Permeability of free space & $\mu_0 = 4\pi\times\SI{e-7}{T\cdot m/A}$ \\
Index of refraction, air & $n \approx 1.00$ \\
Index of refraction, water & $n \approx 1.33$ \\
Index of refraction, glass & $n \approx 1.5$ \\
\end{tabular}

\end{document}
